% GENERATED FILE. Edit canonical lessons, then run npm run generate:books.
% canonical-source-hash: fnv1a64:82f386dd777fd5d0
% canonical-lessons: SA-C274-meet-one-rental-exchange, SA-C274-rental-first-retrieval, SA-C274-rental-final-retrieval

\chapter{One rental exchange, held as a whole}
\label{ch:sa-bounded-rental}
\hlchaptermodality{274}

\begin{chapteropening}
\textbf{By the end of this chapter:} \emph{I can ask and answer one attested question about a rented house.}

Complete chapter 274 by independently retrieving the bounded rental exchange a second time.
\end{chapteropening}

\section[\texorpdfstring{\sk{किं} \sk{तत्} \sk{भाटकितं} … (kiṃ tat …)}{किं तत् भाटकितं … (kiṃ tat …)}]{One rental exchange, held as a whole}

\label{lesson:SA-C274-meet-one-rental-exchange}



\begin{quote}
Recall \textbf{\sk{गृहम्}}, “house” or “home,” and \textbf{\sk{क्रीणामि}}, “I buy.” Then recall the familiar household actions \textbf{\sk{पचति}}, \textbf{\sk{प्रक्षालयति}}, and \textbf{\sk{मार्जयति}}: cooking, washing, and cleaning. This is your known house, buying, and household-work field.
\end{quote}

\subsection*{You'll want to know: one complete attested exchange}
\begin{quote}
\textbf{\sk{किं} \sk{तत्} \sk{भाटकितं} \sk{गृहम्}?}
\end{quote}

\begin{quote}
\emph{kiṃ tat bhāṭakitaṃ gṛham?}
\end{quote}

\begin{quote}
“Is that a rented house?”
\end{quote}

\begin{quote}
\textbf{\sk{न}, \sk{अहं} \sk{भाटकं} \sk{न} \sk{ददामि।}}
\end{quote}

\begin{quote}
\emph{na, ahaṃ bhāṭakaṃ na dadāmi.}
\end{quote}

\begin{quote}
“No, I do not pay rent.”
\end{quote}

The \href{https://sanskritkalp.com/samskrit-vaartaalaap/kim}{direct conversation source} supplies this complete attested exchange. A 1954 \href{https://sanskritdocuments.org/doc\_z\_misc\_major\_works/vyavahAraprashnottarI.html}{\emph{Saṃskṛtavyavahārapraśnottarī} conversation} and its Hindi gloss corroborate the rented-house and rent-payment function.

The older volunteer rendering prints \textbf{\sk{भाटं}}, with an apparent dropped syllable. The Government of India's \href{https://www.cstt.education.gov.in/sites/default/files/fundamintal-administrative-terminology-english-sanskrit.pdf}{Commission for Scientific and Technical Terminology} gives \textbf{\sk{भाटकम्}} for “hire, rent” and \textbf{\sk{गृहभाटकम्}} for “house rent.” The conversation above independently preserves \textbf{\sk{भाटकं}}, so this lesson uses the cross-checked complete exchange rather than copying that apparent omission.

\begin{grammarlens}[title={Grammar Lens: keep the exchange bounded}]
Hear the question and answer as one practical pair. The lesson does not turn \textbf{\sk{भाटकितम्}}, \textbf{\sk{गृहम्}}, or \textbf{\sk{भाटकम्}} into a productive participle or case paradigm, and it does not ask you to derive their sandhi.
\end{grammarlens}

\subsection*{Your turn}
Read the question and answer once slowly. Then say them together as one exchange about whether a house is rented.

\begin{tcolorbox}[breakable,colback=teal!4,colframe=teal!35!black,title={Before you move on}]
Say the complete pair once more. Keep the attested wording together.
\end{tcolorbox}
\section[(independent rental-exchange …]{Retrieve the rental exchange: round one}

\label{lesson:SA-C274-rental-first-retrieval}



\begin{quote}
Prepare to retrieve both turns from their meaning.
\end{quote}

\subsection*{Your turn: independent retrieval}
No Sanskrit model is shown before you answer.

You visit someone at home. Ask, “Is that a rented house?” Then answer, “No, I do not pay rent.” Say the Sanskrit rental exchange.

\begin{grammarlens}[title={What you've built: check after answering}]
Now check: \textbf{\sk{किं} \sk{तत्} \sk{भाटकितं} \sk{गृहम्}? — \sk{न}, \sk{अहं} \sk{भाटकं} \sk{न} \sk{ददामि।}} (\emph{kiṃ tat bhāṭakitaṃ gṛham? — na, ahaṃ bhāṭakaṃ na dadāmi.})
\end{grammarlens}

\begin{tcolorbox}[breakable,colback=teal!4,colframe=teal!35!black,title={Before you move on}]
One successful retrieval is not evidence of productive participle, case, or sandhi mastery, a general renting repertoire, housing competence, human validation, learner attainment, or exam readiness.
\end{tcolorbox}
\section[(independent rental-exchange …]{Retrieve the rental exchange: round two}

\label{lesson:SA-C274-rental-final-retrieval}



\begin{quote}
Retrieve the same two turns in a fresh home context.
\end{quote}

\subsection*{Your turn: independent retrieval in a fresh context}
No Sanskrit model is shown before you answer.

A friend points to another home. Ask, “Is that a rented house?” Then answer, “No, I do not pay rent.” Say the Sanskrit rental exchange.

\begin{grammarlens}[title={What you've built: check after answering}]
Now check: \textbf{\sk{किं} \sk{तत्} \sk{भाटकितं} \sk{गृहम्}? — \sk{न}, \sk{अहं} \sk{भाटकं} \sk{न} \sk{ददामि।}} (\emph{kiṃ tat bhāṭakitaṃ gṛham? — na, ahaṃ bhāṭakaṃ na dadāmi.})
\end{grammarlens}

\begin{tcolorbox}[breakable,colback=teal!4,colframe=teal!35!black,title={Before you move on}]
Two bounded retrievals are not evidence of productive participle, case, or sandhi mastery, a general renting repertoire, housing competence, human validation, learner attainment, or exam readiness.
\end{tcolorbox}
